\documentclass[../../../include/open-logic-section]{subfiles}

\begin{document}

\olfileid[hi]{sth}{ord-arithmetic}{add}
\olsection{क्रमसूचक योग}

मान लें कि हम $\alpha$ और $\beta$ को जोड़ना चाहते हैं। हम $\alpha$ की एक
प्रति के ठीक बाद $\beta$ की एक प्रति रख सकते हैं। (हमें \emph{प्रतियाँ}
लेनी होंगी, क्योंकि \olref[ordinals][basic]{ordinalsaresubsets} से जानते हैं
कि या तो $\alpha \subseteq \beta$ अथवा $\beta \subseteq \alpha$।) इसका सहज
प्रभाव चरणों के एक $\alpha$-अनुक्रम से गुजरना और \emph{फिर} एक
$\beta$-अनुक्रम से गुजरना है। परिणामी अनुक्रम सुव्यवस्थित होगा; अतः
\olref[ordinals][ordtype]{thmOrdinalRepresentation} से वह किसी (अद्वितीय)
क्रमसूचक के समाकृतिक है। उस क्रमसूचक को $\alpha$ और $\beta$ का \emph{योग}
माना जा सकता है।

क्रमसूचक योग के पीछे यही सहज विचार है। उसे कठोर रूप से परिभाषित करने के लिए
हम समुच्चयों की \emph{प्रतियाँ} लेने के विचार से आरंभ करते हैं। यहाँ विचार
यह है कि कौन-सी वस्तु कहाँ से आई इसका पता रखने के लिए मनमाने चिह्न $0$ और
$1$ उपयोग किए जाएँ:

\begin{defn}\ollabel{defdissum}
$A$ और $B$ का \emph{असंयुक्त योग}
$A \disjointsum B = (A\times \{0\}) \cup (B \times \{1\})$ है।
\end{defn}

अब क्रमसूचकों के युग्मों पर एक क्रम परिभाषित करें:

\begin{defn}
किन्हीं क्रमसूचकों $\alpha_1, \alpha_2, \beta_1, \beta_2$ के लिए कहें:
\begin{align*}
	\tuple{\alpha_1, \alpha_2} \rlexless \tuple{\beta_1, \beta_2}
	\text{ तभी और केवल तभी जब }&
	\text{या तो $\alpha_2 \in \beta_2$}\\
	& \text{अथवा $\alpha_2 = \beta_2$ और $\alpha_1 \in \beta_1$, दोनों}
\end{align*} 
\end{defn}

यह \emph{प्रतिलोम शब्दकोशीय} क्रम है, क्योंकि क्रम पहले दूसरे अवयव और फिर
पहले अवयव से किया जाता है। अब याद करें कि हम $\alpha \ordplus \beta$ को
$\alpha$ की प्रति के बाद $\beta$ की प्रति वाले क्रम का क्रम-प्रकार परिभाषित
करना चाहते थे। इसके लिए हम कहते हैं:

\begin{defn}\ollabel{defordplus}
किन्हीं क्रमसूचकों $\alpha$, $\beta$ का योग
$\alpha \ordplus \beta = \ordtype{\alpha \disjointsum \beta,
\rlexless}$ है।
\end{defn}
\noindent
ध्यान दें कि यहाँ हमने संकेतन का थोड़ा दुरुपयोग किया है; कठोर रूप में
``$\rlexless$'' के स्थान पर
``$\Setabs{\tuple{x,y}\in \alpha\disjointsum\beta}{x \rlexless y}$''
लिखना चाहिए। किंतु संक्षेप के लिए आगे भी संकेतन का यह दुरुपयोग करते रहेंगे।

निम्न परिणाम \olref[ordinals][ordtype]{thmOrdinalRepresentation} के साथ
पुष्टि करता है कि हमारी परिभाषा सुगठित है:

\begin{lem}\ollabel{ordsumlessiswo} 
किन्हीं क्रमसूचकों $\alpha$ और $\beta$ के लिए
$\tuple{\alpha \disjointsum \beta, \rlexless}$ सुव्यवस्था है।
\end{lem}

\begin{proof}
स्पष्टतः $\rlexless$, $\alpha \disjointsum \beta$ पर संबद्ध है। उसकी
सुस्थापितता दिखाने के लिए कोई रिक्तेतर
$X \subseteq \alpha \disjointsum \beta$ नियत करें। $Y$ को $X$ का वह
उपसमुच्चय मानें जिसके सदस्यों का दूसरा निर्देशांक यथासंभव छोटा है; अर्थात
$Y = \Setabs{\tuple{\gamma, i} \in X}{(\forall \tuple{\delta, j} \in X)i
\leq j}$। अब $Y$ का लघुतम प्रथम निर्देशांक वाला अवयव चुनें।
\end{proof}
\noindent 
इस प्रकार हमारे पास क्रमसूचक योग की अच्छी, स्पष्ट परिभाषा है। एक
अनाश्चर्यजनक तथ्य यह है (याद करें कि
\olref[z][infinity-again]{defnomega} से $1 = \{0\}$):
\begin{prop}
किसी भी क्रमसूचक $\alpha$ के लिए
$\alpha \ordplus 1 = \ordsucc{\alpha}$।
\end{prop}

\begin{proof}
$\ordsucc{\alpha} = \alpha \cup \{\alpha\}$ से
$\alpha\disjointsum1 = (\alpha \times \{0\}) \disjointsum
(\{0\} \times \{1\})$ में उस समाकृतिकता $f$ पर विचार करें जो
$\gamma \in \alpha$ के लिए $f(\gamma) = \tuple{\gamma, 0}$ और
$f(\alpha) = \tuple{0, 1}$ से दी गई है।
\end{proof}
\noindent
इसके अतिरिक्त यह दिखाना सरल है कि योग कुछ पुनरावर्ती शर्तों का पालन करता है:

\begin{lem}\ollabel{ordadditionrecursion}
किन्हीं क्रमसूचकों $\alpha, \beta$ के लिए हमारे पास है:
\begin{align*}
	\alpha\ordplus 0 &= \alpha\\
	\alpha \ordplus  (\beta\ordplus 1) &= (\alpha \ordplus  \beta) \ordplus  1\\
	\alpha  \ordplus  \beta &= \supstrict_{\delta < \beta}(\alpha \ordplus  \delta) && \text{यदि $\beta$ सीमा क्रमसूचक है}
\end{align*}
\end{lem}

\begin{proof}
हम प्रत्येक मामला जाँचते हैं; पहले:
\begin{align*}
	\alpha \ordplus  0 
	& = \ordtype{(\alpha \times \{0\}) \cup (0 \times \{1\}), \rlexless} \\
	&= \ordtype{(\alpha \times \{0\}) \cup \{0\}, \rlexless}\\
	&= \alpha\\
	\alpha \ordplus (\beta \ordplus  1) 
	%&= \ordtype{(\alpha\times \{0\}) \cup (\ordtype{\beta \ordplus  1}\times \{1\}), \rlexless} \\
	&= \ordtype{(\alpha\times \{0\}) \cup (\ordsucc{\beta}\times \{1\}), \rlexless} \\
	&= \ordtype{(\alpha\times \{0\}) \cup (\beta \times \{1\}), \rlexless} \ordplus 1\\
	&= (\alpha \ordplus  \beta) \ordplus  1
\end{align*}
अब $\beta \neq \emptyset$ को सीमा मानें। यदि $\delta < \beta$, तो
$\delta\ordplus 1 < \beta$ भी; अतः $\alpha \ordplus \delta$,
$\alpha \ordplus \beta$ का उचित आरंभिक खंड है। इसलिए
$\alpha \ordplus \beta$, $X = \Setabs{\alpha \ordplus \delta}{\delta <
\beta}$ का कठोर ऊपरी परिबंध है। इसके अतिरिक्त यदि
$\alpha \leq \gamma < \alpha \ordplus \beta$, तो स्पष्टतः किसी
$\delta < \beta$ के लिए $\gamma = \alpha \ordplus \delta$। अतः
$\alpha \ordplus \beta =
\supstrict_{\delta< \beta}(\alpha\ordplus \delta)$।
\end{proof}

किंतु एक उल्लेखनीय तथ्य यह है। क्रमसूचक योग परिभाषित करने के लिए हम
\emph{इसके बजाय} बस परिमितेतर पुनरावर्तन प्रमेय का उपयोग करके ठीक
\olref{ordadditionrecursion} में दिए पुनरावर्तन समीकरण निर्धारित कर सकते थे
(हालाँकि ``$\beta \ordplus 1$'' के बजाय ``$\ordsucc{\beta}$'' का उपयोग
करते)।

अतः क्रमसूचकों पर संक्रियाएँ परिभाषित करने के दो भिन्न ढंग हैं। किसी
सुव्यवस्थित समुच्चय को स्पष्टतः बनाकर और उसके क्रम-प्रकार पर विचार करके हम
उन्हें \emph{संश्लेषणात्मक} रूप से परिभाषित कर सकते हैं। अथवा केवल पुनरावर्तन
समीकरण निर्धारित करके उन्हें \emph{पुनरावर्ती} रूप से परिभाषित कर सकते हैं।
किंतु सही ढंग से करने पर परिणाम समान है। क्योंकि
\olref[ordinals][ordtype]{thmOrdinalRepresentation} गारंटी देता है कि ये
पुनरावर्तन समीकरण \emph{अद्वितीय} क्रमसूचक निर्धारित करते हैं।

कई अर्थों में क्रमसूचक अंकगणित प्राकृतिक संख्याओं के योग जैसा ही व्यवहार
करता है। उदाहरणतः हम निम्न सिद्ध कर सकते हैं:

\begin{lem}\ollabel{ordinaladditionisnice}
यदि $\alpha, \beta, \gamma$ क्रमसूचक हैं, तो:
\begin{enumerate}
	\item\ollabel{ordaddition1} यदि $\beta < \gamma$, तो
	$\alpha \ordplus \beta < \alpha \ordplus \gamma$;
	\item\ollabel{ordaddition2} यदि
	$\alpha \ordplus \beta = \alpha\ordplus \gamma$, तो $\beta = \gamma$;
	\item\ollabel{ordaddition3}  $\alpha \ordplus  (\beta \ordplus
	\gamma) = (\alpha \ordplus \beta) \ordplus \gamma$; अर्थात योग
	साहचर्यात्मक है;
	\item\ollabel{ordaddition4} यदि $\alpha \leq \beta$, तो
	$\alpha \ordplus \gamma \leq \beta \ordplus \gamma$।
\end{enumerate}
\end{lem}

\begin{proof}
हम \olref{ordaddition3} सिद्ध करते हैं और शेष अभ्यास के लिए छोड़ते हैं।
प्रमाण $\gamma$ पर सरल परिमितेतर आगमन द्वारा है, जिसमें
\olref{ordadditionrecursion} का उपयोग होता है। जब $\gamma = 0$:
\[
(\alpha \ordplus  \beta) \ordplus  0 = \alpha \ordplus  \beta  = \alpha \ordplus  (\beta \ordplus  0)
\]
जब $\gamma = \delta\ordplus 1$, आगमन के लिए मान लें कि
$(\alpha \ordplus \beta) \ordplus \delta
= \alpha \ordplus (\beta \ordplus \delta)$; अब
\olref{ordadditionrecursion} का तीन बार उपयोग करके:
\begin{align*}
	(\alpha \ordplus  \beta) \ordplus  (\delta \ordplus  1) & = ((\alpha \ordplus  \beta) \ordplus  \delta)\ordplus 1\\
	& = (\alpha \ordplus  (\beta \ordplus  \delta)) \ordplus  1\\
	& = \alpha \ordplus  ((\beta \ordplus  \delta)\ordplus 1)\\
	& = \alpha \ordplus  (\beta \ordplus  (\delta\ordplus 1))
\end{align*}	
जब $\gamma$ सीमा क्रमसूचक है, आगमन के लिए मान लें कि यदि
$\delta \in \gamma$, तो
$(\alpha \ordplus \beta) \ordplus \delta
= \alpha \ordplus (\beta \ordplus \delta)$; अब:
\begin{align*}
	(\alpha \ordplus  \beta) \ordplus  \gamma & = \supstrict_{\delta < \gamma}((\alpha \ordplus  \beta) \ordplus  \delta) \\
	&= \supstrict_{\delta < \gamma}(\alpha \ordplus  (\beta \ordplus  \delta))\\
	&= \alpha \ordplus  \supstrict_{\delta < \gamma}(\beta \ordplus  \delta)\\
	& = \alpha \ordplus  (\beta \ordplus  \gamma)
\end{align*}
\end{proof}

\begin{prob}
\olref[sth][ord-arithmetic][add]{ordinaladditionisnice} का शेष सिद्ध कीजिए।
\end{prob}

इन अर्थों में क्रमसूचक योग बहुत परिचित होना चाहिए। किंतु एक निर्णायक अर्थ
में क्रमसूचक योग प्राकृतिक संख्याओं के योग जैसा \emph{नहीं} है।

\begin{prop}\ollabel{ordsumnotcommute}
क्रमसूचक योग क्रमविनिमेय {नहीं} है;
$1 \ordplus \omega = \omega < \omega \ordplus 1$।
\end{prop}

\begin{proof}
ध्यान दें कि
$1 \ordplus \omega = \supstrict_{n < \omega} (1 \ordplus n)
= \omega \in \omega \cup \{\omega\} = \ordsucc{\omega}
= \omega \ordplus 1$।
\end{proof}
\noindent 
यह आरंभ में आश्चर्यजनक लगे, किंतु \emph{लगना नहीं चाहिए}। एक ओर
$1 \ordplus \omega$ पर विचार करते समय आप सभी प्राकृतिक संख्याओं के
\emph{पहले} एक अतिरिक्त अवयव रखने से मिले क्रम-प्रकार के विषय में सोचते हैं।
\olref[sfr][infinite][hilbert]{sec} में हिल्बर्ट के होटल के साथ किए तर्क की
तरह सहजतः इस अतिरिक्त प्रथम अवयव से समग्र क्रम-प्रकार पर कोई प्रभाव नहीं
पड़ना चाहिए। दूसरी ओर $\omega \ordplus 1$ पर विचार करते समय आप सभी
प्राकृतिक संख्याओं के \emph{बाद} एक अतिरिक्त अवयव रखने से मिले क्रम-प्रकार
के विषय में सोचते हैं। और वह मूलतः भिन्न वस्तु है!{}

\end{document}
